\documentclass[10pt,a4paper]{amsart}
\usepackage{amsmath,amssymb,amsthm,graphicx}
% Retained preamble for the Stacks Project chapter snapshot in this directory.
% The chapter's own preamble.tex is a copy of the shared Stacks Project preamble; the block
% below is sliced MECHANICALLY, line for line, from it. Nothing was added, removed or
% redefined; the wrapper class, the xr-hyper/hyperref wiring and the externaldocument list
% are omitted because the assembler supplies the document class.
%
% lines taken from preamble.tex: 3, 5, 9, 11, 12, 156, 157, 158, ... (51 lines total)
%
% package coverage restored from the chapter's own preamble (2026-10-13): xy, xy-2cell

\usepackage[all]{xy}
\xyoption{2cell}
\UseAllTwocells

\newenvironment{reference}{\comment}{\endcomment}
\newenvironment{slogan}{\comment}{\endcomment}
\newenvironment{history}{\comment}{\endcomment}
\theoremstyle{plain}
\newtheorem{theorem}[subsection]{Theorem}
\newtheorem{proposition}[subsection]{Proposition}
\newtheorem{lemma}[subsection]{Lemma}
\theoremstyle{definition}
\newtheorem{definition}[subsection]{Definition}
\newtheorem{example}[subsection]{Example}
\newtheorem{exercise}[subsection]{Exercise}
\newtheorem{situation}[subsection]{Situation}
\theoremstyle{remark}
\newtheorem{remark}[subsection]{Remark}
\newtheorem{remarks}[subsection]{Remarks}
\numberwithin{equation}{subsection}
\def\lim{\mathop{\mathrm{lim}}\nolimits}
\def\colim{\mathop{\mathrm{colim}}\nolimits}
\def\Spec{\mathop{\mathrm{Spec}}}
\def\Hom{\mathop{\mathrm{Hom}}\nolimits}
\def\Ext{\mathop{\mathrm{Ext}}\nolimits}
\def\SheafHom{\mathop{\mathcal{H}\!\mathit{om}}\nolimits}
\def\SheafExt{\mathop{\mathcal{E}\!\mathit{xt}}\nolimits}
\def\Sch{\mathit{Sch}}
\def\Mor{\mathop{\mathrm{Mor}}\nolimits}
\def\Ob{\mathop{\mathrm{Ob}}\nolimits}
\def\Sh{\mathop{\mathit{Sh}}\nolimits}
\def\NL{\mathop{N\!L}\nolimits}
\def\CH{\mathop{\mathrm{CH}}\nolimits}
\def\proetale{{pro\text{-}\acute{e}tale}}
\def\etale{{\acute{e}tale}}
\def\QCoh{\mathit{QCoh}}
\def\Ker{\mathop{\mathrm{Ker}}}
\def\Im{\mathop{\mathrm{Im}}}
\def\Coker{\mathop{\mathrm{Coker}}}
\def\Coim{\mathop{\mathrm{Coim}}}
\DeclareMathSymbol{\boxtimes}{\mathbin}{AMSa}{"02}
\def\QCohstack{\mathcal{QC}\!\mathit{oh}}
\def\Cohstack{\mathcal{C}\!\mathit{oh}}
\def\Spacesstack{\mathcal{S}\!\mathit{paces}}
\def\Quotfunctor{\mathrm{Quot}}
\def\Hilbfunctor{\mathrm{Hilb}}
\def\Curvesstack{\mathcal{C}\!\mathit{urves}}
\def\Polarizedstack{\mathcal{P}\!\mathit{olarized}}
\def\Complexesstack{\mathcal{C}\!\mathit{omplexes}}
\def\Pic{\mathop{\mathrm{Pic}}\nolimits}
\def\Picardstack{\mathcal{P}\!\mathit{ic}}
\def\Picardfunctor{\mathrm{Pic}}
\def\Deformationcategory{\mathcal{D}\!\mathit{ef}}
\title{Four properties of module maps that descend to a finite stage of a colimit}
\author{The Stacks Project Authors}
\date{stacks-project:algebra.tex:local-snapshot-20261012; GNU FDL 1.2 or later}
\begin{document}
\maketitle
\noindent\textit{Editorial note.} All delivered mathematical text is the author's own, copied from the pinned Stacks Project chapter file algebra.tex; no statement or proof step was rewritten, reordered or removed.\par\medskip
\noindent\textit{Editorial note.} Preamble provenance, as required for this source: the frozen snapshot stages the chapter's own preamble file beside it (retained-preamble.tex, 63 lines, 8116 bytes, the same shared Stacks-wide preamble the published chow and etale-cohomology cards use); the first line of the pinned chapter file is \textbackslash{}input\{preamble\}, so no preamble text lives inside the chapter itself. The card takes that file's lines 8-18 (the verbatim package and the three \textbackslash{}comment-based environments reference/slogan/history at lines 9-12, section 13.2, together with the Xy-pic setup) and its lines 154 to the runtime end of the file (the native theorem declarations and mathematical macros). Not one character was changed, no package was added, no macro was defined, and the excerpt serves only as preamble, never as a span coordinate.\par\medskip
\noindent\textit{Editorial note.} The same-chapter pointers are located at their \textbackslash{}label lines inside the pinned chapter file; the cross-chapter pointers are located in the project's own tag list staged in the snapshot (sources/stacks-chow/tags.txt), at the verbatim lines carrying their canonical tags.\par\medskip
\noindent\textit{Editorial note.} The counted spans contain no \textbackslash{}ref and no \textbackslash{}cite, so the card carries no navigation stub and no bibliography.\par\medskip
\noindent\textit{Allowed standard prerequisites (editorial).}\begin{itemize}
\item Finitely many generators spread out along a colimit. Used in part (1) of the counted proof (lines 32233-32237). Moves the finitely many equations defining a vanishing of two maps into one stage.
\item A finite presentation of the target to build a retraction. Used in part (3) and (4) of the counted proof. Produces the map on a stage from a map on the colimit.
\end{itemize}\medskip
\begin{lemma}
\label{lemma-module-map-property-in-colimit}
Let $A$ be a ring and let $M, N$ be $A$-modules.
Suppose that $R = \colim_{i \in I} R_i$ is a directed colimit
of $A$-algebras.
\begin{enumerate}
\item If $M$ is a finite $A$-module, and $u, u' : M \to N$ are
$A$-module maps such that
$u \otimes 1 = u' \otimes 1 : M \otimes_A R \to N \otimes_A R$
then for some $i$ we have
$u \otimes 1 = u' \otimes 1 : M \otimes_A R_i \to N \otimes_A R_i$.
\item If $N$ is a finite $A$-module and $u : M \to N$ is an $A$-module
map such that $u \otimes 1 : M \otimes_A R \to N \otimes_A R$ is surjective,
then for some $i$ the map $u \otimes 1 : M \otimes_A R_i \to N \otimes_A R_i$
is surjective.
\item If $N$ is a finitely presented $A$-module, and
$v : N \otimes_A R \to M \otimes_A R$ is an $R$-module
map, then there exists an $i$ and an $R_i$-module map
$v_i : N \otimes_A R_i \to M \otimes_A R_i$ such that $v = v_i \otimes 1$.
\item If $M$ is a finite $A$-module, $N$ is a finitely presented $A$-module,
and $u : M \to N$ is an $A$-module map such that
$u \otimes 1 : M \otimes_A R \to N \otimes_A R$ is an isomorphism, then
for some $i$ the map $u \otimes 1 : M \otimes_A R_i \to N \otimes_A R_i$
is an isomorphism.
\end{enumerate}
\end{lemma}

\begin{proof}
To prove (1) assume $u$ is as in (1) and
let $x_1, \ldots, x_m \in M$ be generators. Since
$N \otimes_A R = \colim_i N \otimes_A R_i$
we may pick an $i \in I$ such that $u(x_j) \otimes 1 = u'(x_j) \otimes 1$
in $M \otimes_A R_i$, $j = 1, \ldots, m$.
For such an $i$ we have
$u \otimes 1 = u' \otimes 1 : M \otimes_A R_i \to N \otimes_A R_i$.

\medskip\noindent
To prove (2) assume $u \otimes 1$ surjective and
let $y_1, \ldots, y_m \in N$ be generators. Since
$N \otimes_A R = \colim_i N \otimes_A R_i$
we may pick an $i \in I$ and $z_j \in M \otimes_A R_i$, $j = 1, \ldots, m$
whose images in $N \otimes_A R$ equal $y_j \otimes 1$.
For such an $i$ the map $u \otimes 1 : M \otimes_A R_i \to N \otimes_A R_i$
is surjective.

\medskip\noindent
To prove (3) let $y_1, \ldots, y_m \in N$ be generators. Let
$K = \Ker(A^{\oplus m} \to N)$ where the map is given by
the rule $(a_1, \ldots, a_m) \mapsto \sum a_j x_j$. Let $k_1, \ldots, k_t$
be generators for $K$. Say $k_s = (k_{s1}, \ldots, k_{sm})$.
Since $M \otimes_A R = \colim_i M \otimes_A R_i$
we may pick an $i \in I$ and $z_j \in M \otimes_A R_i$, $j = 1, \ldots, m$
whose images in $M \otimes_A R$ equal $v(y_j \otimes 1)$.
We want to use the $z_j$ to define the map
$v_i : N \otimes_A R_i \to M \otimes_A R_i$.
Since $K \otimes_A R_i \to R_i^{\oplus m} \to N \otimes_A R_i \to 0$
is a presentation, it suffices to check that $\xi_s = \sum_j k_{sj}z_j$ is
zero in $M \otimes_A R_i$ for each $s = 1, \ldots, t$. This may not
be the case, but since the image of $\xi_s$ in $M \otimes_A R$ is zero
we see that it will be the case after increasing $i$ a bit.

\medskip\noindent
To prove (4) assume $u \otimes 1$ is an isomorphism, that
$M$ is finite, and that $N$ is finitely presented.
Let $v : N \otimes_A R \to M \otimes_A R$ be an inverse to
$u \otimes 1$. Apply part (3) to get a map
$v_i : N \otimes_A R_i \to M \otimes_A R_i$ for some $i$.
Apply part (1) to see that, after increasing $i$ we have
$v_i \circ (u \otimes 1) = \text{id}_{M \otimes_R R_i}$ and
$(u \otimes 1) \circ v_i = \text{id}_{N \otimes_R R_i}$.
\end{proof}

\end{document}
